\section{استخر ریسه و ذخیره‌ی محلی ریسه}

\subsection*{بخش آ – سرور با استخر ۴ ریسه‌ای}

صف \lr{FIFO} است و هر درخواست تا پایان دو \code{part} خود، یک ریسه را در \code{pool\_wait} مسدود نگه می‌دارد؛ \code{part}ها به انتهای صف می‌روند.
\begin{itemize}
  \item \textbf{۳ درخواست: همه تمام می‌شوند.} سه ریسه درخواست‌ها را می‌گیرند و مسدود می‌شوند؛ ریسه‌ی چهارم آزاد است و ۶ \code{part} را یکی‌یکی اجرا می‌کند.
  \item \textbf{۸ درخواست: بن‌بست؛ هیچ درخواستی تمام نمی‌شود.} چهار ریسه \lr{R1} تا \lr{R4} را برمی‌دارند و همه مسدود می‌شوند. \code{part}ها پشت \lr{R5} تا \lr{R8} در صف‌اند و ریسه‌ی آزادی نیست.
\end{itemize}
(با ۴ درخواست یا بیشتر همین رخ می‌دهد؛ حد امن ۳ است.)

\subsection*{بخش ب – \lr{part} هم دو زیرکار می‌سازد}

\textbf{دست‌کم ۴ ریسه.} درخواست و دو \code{part} هر کدام منتظر زیرکارهایشان می‌مانند، یعنی تا ۳ ریسه هم‌زمان مسدودند و برای تکه‌های محاسباتی دست‌کم یک ریسه‌ی آزاد دیگر لازم است.
\begin{itemize}
  \item با ۳ ریسه بن‌بست قطعی است: صف \lr{FIFO} دو \code{part} را پیش از تکه‌های ریز به دو ریسه‌ی آزاد می‌دهد؛ هر سه مسدود می‌شوند و ۴ تکه‌ی محاسباتی ریسه ندارند.
  \item با ۴ ریسه همیشه یکی آزاد است. تکه‌های ریز منتظر چیزی نیستند و تمام می‌شوند، پس همه‌چیز تمام می‌شود.
\end{itemize}

\subsection*{بخش ج – متغیر محلی ریسه}

\code{a} با \code{\_\_thread} برای هر ریسه (از جمله \lr{main}) نسخه‌ی جدا دارد که از ۰ شروع می‌شود؛ \code{b} سراسری و مشترک است. \code{join}ها اجرا را پشت‌سرهم می‌کنند:
\begin{center}\small\atlrows
\begin{tabular}{@{}cc>{\raggedleft\arraybackslash}p{3.4cm}@{}}
  \rowcolor{atl-primary}\atlhead{فراخوانی} & \atlhead{خروجی} & \atlhead{توضیح}\\
  \lr{t1} & \lr{a=1 b=1} & \lr{a} ریسه‌ی \lr{t1}\\
  \lr{t2} & \lr{a=1 b=2} & \lr{a} تازه، \lr{b} ادامه\\
  \lr{main}، بار اول & \lr{a=1 b=3} & \lr{a} ریسه‌ی \lr{main}\\
  \lr{main}، بار دوم & \lr{a=2 b=4} & همان \lr{a} افزایش می‌یابد\\
\end{tabular}
\end{center}

\subsection*{بخش د – مشکل \lr{errno} سراسری}

اگر \code{errno} متغیر سراسری معمولی بود، ریسه‌ها یکدیگر را خراب می‌کردند. ترتیب اجرا:
\begin{enumerate}
  \item \lr{T1}: \code{open} شکست می‌خورد و \code{errno = ENOENT} می‌شود.
  \item تعویض زمینه؛ \lr{T2}: فراخوانی دیگری شکست می‌خورد و \code{errno = EACCES} می‌شود.
  \item \lr{T1} ادامه می‌دهد و \code{errno} را می‌خواند: \code{EACCES} نه \code{ENOENT}؛ علت خطا اشتباه تشخیص داده می‌شود.
\end{enumerate}
راه‌حل (ایده‌ی بخش ج): \code{errno} را متغیر محلی ریسه کنیم تا هر ریسه نسخه‌ی خودش را داشته باشد. \lr{glibc} نیز همین کار را می‌کند.
